\subsection{Differentiation and Basic Theorems}

% **Role:** calculus as local linearization plus global geometric integration.
%
% Derivatives, MVT, Taylor theorem, Riemann integration, FTC, improper integrals.

One-variable calculus studies functions on intervals through derivatives, mean-value theorems, Riemann integration, and the FTC, then points toward vector-valued estimates and linear approximation.

\begin{definition} \label{def:derivative-real-function}\label{def:one-variable-differentiation-core-source-todo} \hlt{Derivative of a Real Function}\\
Let $f$ be real-valued on $[a,b]$ and let $x\in[a,b]$. For $a<t<b$ and $t\neq x$, set $\phi_x(t)\defeq\frac{f(t)-f(x)}{t-x}$.\\
If $\lim_{t\to x}\phi_x(t)$ exists, denote \hlt{derivative} by $f'(x)$; then $f$ is \hlt{differentiable at $x$}.\\
If this holds for every $x\in E\subseteq[a,b]$, then $f$ is \hlt{differentiable on $E$}.
\end{definition}

\begin{figure}[H]
  \centering
  \begin{tikzpicture}[
  scale=0.92,
  >=Stealth,
  every node/.style={font=\scriptsize},
  axis/.style={->, gray!70, line width=0.45pt},
  domainline/.style={blue!65!black, line width=1.0pt},
  approach/.style={->, orange!80!black, line width=0.7pt},
  point/.style={circle, fill=black, inner sep=1.15pt},
  guide/.style={gray!55, dashed, line width=0.45pt},
  labelbox/.style={fill=white, fill opacity=0.9, text opacity=1, inner sep=1.2pt}
]
  \begin{scope}
    \draw[axis] (0,0) -- (4.4,0) node[right] {$\R$};
    \foreach \x in {0.75,2.2,3.6} {
      \node[point] at (\x,0) {};
    }
    \draw[guide] (1.62,-0.26) rectangle (2.78,0.26);
    \node[above=2pt] at (2.2,0) {$x_0$};
  \end{scope}

  \begin{scope}[xshift=5.3cm]
    \draw[axis] (0,0) -- (4.4,0) node[right] {$\R$};
    \draw[domainline] (0.75,0) -- (4.05,0);
    \node[point] at (0.75,0) {};
    \node[above=2pt] at (0.75,0) {$x_0$};
    \node[labelbox, above=11pt] at (2.45,0.26) {$X=[x_0,\infty)$};
    \draw[approach] (3.55,0.32) -- (2.75,0.07);
    \draw[approach] (2.7,0.32) -- (1.95,0.07);
    \draw[approach] (1.9,0.32) -- (1.05,0.07);
  \end{scope}
\end{tikzpicture}

  \caption{Domain matters: isolated points give no punctured approach; $[x_0,\infty)$ gives a right-hand one.}
  \label{fig:one-variable-domain-restriction}
\end{figure}

\begin{perspective} \label{per:one-variable-domain-of-differentiability} \hlt{Derivative Depends on the Domain}\\
For $f:X\to\R$, $f'(x_0)$ is the limit of $\frac{f(x)-f(x_0)}{x-x_0}$ as $x\to x_0$ in $X\setminus\{x_0\}$; see Definition \ref{def:function-limit-metric}. Hence $x_0$ must be a limit point of $X$: isolated points give no derivative, while half-interval domains give one-sided derivatives.
\end{perspective}

\begin{figure}[H]
  \centering
  \begin{tikzpicture}[
  scale=0.37,
  >=Stealth,
  every node/.style={font=\scriptsize},
  axis/.style={->, gray!70, line width=0.45pt},
  curve/.style={blue!70!black, line width=0.8pt},
  secant/.style={orange!80!black, line width=0.7pt},
  tangent/.style={green!45!black, line width=0.7pt},
  guide/.style={gray!45, dashed, line width=0.4pt},
  point/.style={circle, fill=gray!85, inner sep=1.4pt},
  labelbox/.style={fill=white, fill opacity=0.88, text opacity=1, inner sep=1.2pt}
]
  \draw[axis] (-0.625,0) -- (12.4,0) node[right] {$t$};
  \draw[axis] (0,-0.5) -- (0,8.7) node[above] {$y$};

  \draw[curve, domain=0.25:9.5, samples=90]
    plot (\x,{0.875*(\x/2.5-1.2)*(\x/2.5-1.2)+1.75});

  \coordinate (X) at (3.5,1.785);
  \coordinate (T) at (8.0,5.25);

  \draw[secant] (1.875,0.525) -- (9.3,6.25);
  \draw[tangent] (1.125,1.4525) -- (9.875,2.6775);

  \draw[guide] (3.5,0) -- (X);
  \draw[guide] (8.0,0) -- (T);
  \node[below] at (3.5,0) {$x$};
  \node[below] at (8.0,0) {$t$};

  \node[point] at (X) {};
  \node[point] at (T) {};
  \node[labelbox, anchor=south east] at (3.5,2.45)
    {\textcolor{blue!70!black}{$(x,f(x))$}};
  \node[labelbox, anchor=west] at (8.55,5.05)
    {\textcolor{blue!70!black}{$(t,f(t))$}};

  \node[labelbox, align=center] at (5.75,5.25)
    {\textcolor{orange!75!black}{$\frac{f(t)-f(x)}{t-x}$}};
  \node[labelbox] at (8.65,3.2)
    {\textcolor{green!45!black}{$f'(x)$}};
\end{tikzpicture}

  \caption{The derivative is the limiting slope of secant lines through $(x,f(x))$.}
  \label{fig:derivative-secant-to-tangent}
\end{figure}

\begin{proposition} \label{prop:one-variable-first-order-approximation} \hlt{First-Order Approximation Form}\\
Let $X\subseteq\R$, let $x_0\in X$ be a limit point, let $f:X\to\R$, and let $L\in\R$. Then $f'(x_0)=L$ relative to $X$ iff $\forall \varepsilon>0$ $\exists \delta>0$ such that $\abs{f(x)-f(x_0)-L(x-x_0)}\leq \varepsilon\abs{x-x_0}$ whenever $x\in X$ and $\abs{x-x_0}<\delta$.
\begin{proof}
For $x\neq x_0$, $\frac{\abs{f(x)-f(x_0)-L(x-x_0)}}{\abs{x-x_0}} = \abs{\frac{f(x)-f(x_0)}{x-x_0}-L}$, while $x=x_0$ makes approx. estimate automatic.
\begin{enumerate}[label=\roman*.]
\setlength{\itemsep}{0pt}
\item[$(\Rightarrow)$] If $f'(x_0)=L$, then the displayed right-hand side is at most $\varepsilon$ for all $x\in X\setminus\{x_0\}$ sufficiently near $x_0$. Multiplying by $\abs{x-x_0}$ gives the approximation estimate.
\item[$(\Leftarrow)$] If the approximation estimate holds, then division by $\abs{x-x_0}$ gives $\abs{\frac{f(x)-f(x_0)}{x-x_0}-L}\leq\varepsilon$ for all $x\in X\setminus\{x_0\}$ sufficiently near $x_0$. Hence $f'(x_0)=L$.
\end{enumerate}
\hlt{Structural Commentary}: Differentiability is first-order linear approximation.\\
\hlt{Key Insight}: Multiply or divide by $\abs{x-x_0}$.
\end{proof}
\end{proposition}

\begin{theorem} \label{thm:differentiable-implies-continuous}\label{thm:one-variable-differentiation-main-source-todo} \hlt{Differentiability Implies Continuity}\\
If $f$ is differentiable at $x$, then $f$ is continuous at $x$.
\begin{proof}
For $t\neq x$, $f(t)-f(x)=\frac{f(t)-f(x)}{t-x}(t-x)$. As $t\to x$, the difference quotient tends to $f'(x)$ and $t-x\to0$.\\
By Theorem \ref{thm:function-limit-algebra}, $f(t)-f(x)\to0$, so $f(t)\to f(x)$. Theorem \ref{thm:continuity-as-limit-at-limit-point} gives continuity at $x$.\\
\hlt{Structural Commentary}: A differentiable function has first-order control near the point.\\
\hlt{Key Insight}: The quotient stays near a finite number while $t-x$ shrinks to zero.
\end{proof}
\end{theorem}

\begin{theorem} \label{thm:derivative-algebra-rules} \hlt{Algebra Rules for Derivatives}\\
If $f$ and $g$ are differentiable at $x$, then so are $f+g$ and $fg$; if $g(x)\neq0$, then $f/g$ is differentiable at $x$. Moreover,
\begin{align}
(f+g)'(x)&=f'(x)+g'(x), \nonumber\\
(fg)'(x)&=f'(x)g(x)+f(x)g'(x), \nonumber\\
\left(\frac{f}{g}\right)'(x)&=\frac{g(x)f'(x)-g'(x)f(x)}{g(x)^2}. \nonumber
\end{align}
\begin{proof}
The sum rule is limit algebra. For $h=fg$, $h(t)-h(x)=f(t)\bigl(g(t)-g(x)\bigr)+g(x)\bigl(f(t)-f(x)\bigr)$.\\
Dividing by $t-x$ and letting $t\to x$ gives the product rule, using Theorem \ref{thm:differentiable-implies-continuous} for $f(t)\to f(x)$.\\
For $h=f/g$, continuity of $g$ makes $g(t)\neq0$ near $x$, and $\frac{h(t)-h(x)}{t-x} =\frac{1}{g(t)g(x)} \left[ g(x)\frac{f(t)-f(x)}{t-x} -f(x)\frac{g(t)-g(x)}{t-x} \right]$.\\
Letting $t\to x$ gives the quotient rule.\\
\hlt{Structural Commentary}: Product and quotient rules are difference-quotient identities plus continuity.\\
\hlt{Key Insight}: Difference-quotient identities plus continuity give the algebra rules.
\end{proof}
\end{theorem}

\begin{theorem} \label{thm:chain-rule-one-variable} \hlt{Chain Rule}\\
Let $f$ be continuous on $[a,b]$ with $f'(x)$ existing at $x\in[a,b]$, and let $g$ be defined on an interval containing $f([a,b])$ and differentiable at $f(x)$. If $h(t)=g(f(t))$, then $h'(x)=g'(f(x))f'(x)$.
\begin{proof}
Let $y=f(x)$. For some $u(t)\to0$ and $v(s)\to0$,
\begin{align}
f(t)-f(x)&=(t-x)\bigl(f'(x)+u(t)\bigr), \nonumber\\
g(s)-g(y)&=(s-y)\bigl(g'(y)+v(s)\bigr). \nonumber
\end{align}
With $s=f(t)$, continuity gives $s\to y$, and for $t\neq x$, $\frac{h(t)-h(x)}{t-x} =\bigl(g'(y)+v(f(t))\bigr)\bigl(f'(x)+u(t)\bigr)\to g'(y)f'(x)$.\\
\hlt{Structural Commentary}: Each differentiability statement is a linear approximation plus a vanishing error.\\
\hlt{Key Insight}: Composition multiplies the two first-order parts; the error terms disappear.
\end{proof}
\end{theorem}

\begin{example} \label{ex:oscillatory-derivative-examples} \hlt{Oscillation at a Point}\\
Assuming the usual derivative of $\sin x$, define
\begin{align}
f_1(x)=
\begin{cases}
x\sin(1/x),&x\neq0,\\
0,&x=0,
\end{cases}
\qquad
f_2(x)=
\begin{cases}
x^2\sin(1/x),&x\neq0,\\
0,&x=0.
\end{cases}
\nonumber
\end{align}
For $f_1$, the quotient at $0$ is $\sin(1/t)$, so $f_1'(0)$ does not exist. For $f_2$, $\abs{(f_2(t)-f_2(0))/t}\leq\abs{t}$, so $f_2'(0)=0$.\\
However, for $x\neq0$, $f_2'(x)=2x\sin(1/x)-\cos(1/x)$, so $f_2'$ is not continuous at $0$.
\end{example}

\begin{definition} \label{def:local-extrema} \hlt{Local Extrema}\\
Let $f$ be a real-valued function on a metric space $X$. The function has a \hlt{local maximum} at $p\in X$ if there is $\delta>0$ such that $f(q)\leq f(p)$ whenever $d(p,q)<\delta$. Local minima are defined by reversing the inequality.
\end{definition}

\begin{remark} \label{rem:local-versus-global-extrema} \hlt{Local and Global Extrema}\\
Global extrema are local, but local extrema depend on domain: isolated points are automatic extrema. Derivative tests therefore need an interior interval point; compactness supplies global extrema (Theorem \ref{thm:extreme-value-theorem-compact}).
\end{remark}

\begin{theorem} \label{thm:interior-extremum-derivative-zero} \hlt{Interior Extremum Has Zero Derivative}\\
If $f$ has a local extremum at $x\in(a,b)$ and $f'(x)$ exists, then $f'(x)=0$.
\begin{proof}
For a local maximum, $f(t)\leq f(x)$ near $x$.\\
Thus the quotient $(f(t)-f(x))/(t-x)$ is nonnegative for $t<x$ and nonpositive for $t>x$. Since the two-sided limit exists, $f'(x)\geq0$ and $f'(x)\leq0$, so $f'(x)=0$. For a local minimum, apply the local-maximum case to $-f$.\\
Since $(-f)'(x)=-f'(x)$, we get $-f'(x)=0$, hence $f'(x)=0$.\\
\hlt{Structural Commentary}: Interior extrema force opposite one-sided quotient signs.\\
\hlt{Key Insight}: A two-sided derivative compatible with both signs must be zero.
\end{proof}
\end{theorem}

\begin{theorem} \label{thm:rolles-theorem-one-variable} \hlt{Rolle's Theorem}\\
If $f:[a,b]\to\R$ continuous on $[a,b]$, differentiable on $(a,b)$, and $f(a)=f(b)$, then $\exists x\in(a,b)$ s.t. $f'(x)=0$.
\begin{proof}
If $f$ is constant, any $x\in(a,b)$ works. Otherwise, by Theorem \ref{thm:extreme-value-theorem-compact}, $f$ attains a nonendpoint extremum at some $x\in(a,b)$. Theorem \ref{thm:interior-extremum-derivative-zero} gives $f'(x)=0$.\\
\hlt{Structural Commentary}: Equal endpoints force any nonconstant height change to turn around inside.\\
\hlt{Key Insight}: Compactness gives the extremum; differentiability makes the interior extremum stationary.
\end{proof}
\end{theorem}

\begin{theorem} \label{thm:cauchy-mean-value-theorem} \hlt{Cauchy Mean Value Theorem}\\
Let $f,g$ be continuous on $[a,b]$, differentiable on $(a,b)$. $\exists$ $x\in(a,b)$, $\bigl(f(b)-f(a)\bigr)g'(x)=\bigl(g(b)-g(a)\bigr)f'(x)$.
\begin{proof}
Set $h(t)\defeq\bigl(f(b)-f(a)\bigr)g(t)-\bigl(g(b)-g(a)\bigr)f(t)$.\\
Then $h(a)=h(b)$, $h$ is continuous on $[a,b]$, and $h$ is differentiable on $(a,b)$. By Theorem \ref{thm:rolles-theorem-one-variable}, $h'(x)=0$ for some $x\in(a,b)$. Since $h'(x)=\bigl(f(b)-f(a)\bigr)g'(x)-\bigl(g(b)-g(a)\bigr)f'(x)$, the identity follows.\\
\hlt{Structural Commentary}: Build one auxiliary function whose endpoint values agree.\\
\hlt{Key Insight}: Equal endpoints force an interior horizontal tangent.
\end{proof}
\end{theorem}

\begin{figure}[H]
  \centering
  \begin{tikzpicture}[
  scale=0.72,
  >=Stealth,
  every node/.style={font=\scriptsize},
  axis/.style={->, black!70, line width=0.5pt},
  curve/.style={black, line width=0.85pt},
  secant/.style={black, dashed, dash pattern=on 8pt off 6pt, line width=0.65pt},
  tangent/.style={<->, green!45!black, line width=0.75pt},
  guide/.style={black!55, dashed, line width=0.45pt},
  tick/.style={black!70, line width=0.45pt},
  point/.style={circle, fill=black, inner sep=1.45pt},
  labelbox/.style={fill=white, fill opacity=0.94, text opacity=1, inner sep=1.2pt}
]
  \draw[axis] (-0.55,0) -- (7.25,0);
  \draw[axis] (0,-0.5) -- (0,4.55);

  \coordinate (A) at (1.25,1.05);
  \coordinate (C) at (3.85,1.38);
  \coordinate (B) at (5.65,3.35);

  \draw[curve, <-]
    (0.42,1.42) .. controls (0.72,1.28) and (0.98,1.12) .. (A);
  \draw[curve]
    (A) .. controls (1.80,0.78) and (2.95,0.91) .. (C);
  \draw[curve]
    (C) .. controls (4.75,1.85) and (4.95,2.85) .. (B);
  \draw[curve, ->]
    (B) .. controls (6.05,4.10) and (6.35,4.32) .. (6.90,4.52);

  \draw[secant] (0.62,0.72) -- (6.72,3.91);
  \draw[tangent] (2.82,0.84) -- (5.10,2.03);

  \draw[guide] (C) -- (3.85,0);

  \foreach \x/\lab in {1.25/a,3.85/c,5.65/b} {
    \draw[tick] (\x,0.12) -- (\x,-0.12);
    \node[below] at (\x,-0.12) {$\lab$};
  }

  \node[point] at (A) {};
  \node[point] at (B) {};

  \node[labelbox, anchor=south east] at ($(A)+(-0.06,0.12)$) {$(a,f(a))$};
  \node[labelbox, anchor=south west] at ($(B)+(-0.02,0.08)$) {$(b,f(b))$};
  \node[labelbox, anchor=west] at (5.08,1.76)
    {\textcolor{green!45!black}{$f'(c)=\dfrac{f(b)-f(a)}{b-a}$}};
\end{tikzpicture}

  \caption{The Mean Value Theorem finds an interior tangent slope equal to the endpoint secant slope.}
  \label{fig:mean-value-secant-tangent}
\end{figure}

\begin{corollary} \label{cor:mean-value-theorem} \hlt{Mean Value Theorem}\\
If $f$ is continuous on $[a,b]$, differentiable in $(a,b)$, then $\exists$ $x\in(a,b)$ such that $f(b)-f(a)=(b-a)f'(x)$.
\begin{proof}
Let $g(t)=t$. Then $g'(x)=1$, $g(b)-g(a)=b-a$.\\
Hence Theorem \ref{thm:cauchy-mean-value-theorem} gives $f(b)-f(a)=(b-a)f'(x)$ for some $x\in(a,b)$.\\
\hlt{Structural Commentary}: This is Cauchy's theorem with $g(t)=t$.\\
\hlt{Key Insight}: A secant slope is achieved as a tangent slope somewhere inside.
\end{proof}
\end{corollary}

\begin{corollary} \label{cor:bounded-derivative-lipschitz-one-variable} \hlt{Bounded Derivative Gives Lipschitz Control}\\
If $f:[a,b]\to\R$ is continuous on $[a,b]$, differentiable on $(a,b)$, and $\abs{f'}\leq M$ on $(a,b)$, then $f$ is $M$-Lipschitz on $[a,b]$, hence uniformly continuous.
\begin{proof}
For $x<y$, Corollary \ref{cor:mean-value-theorem} gives $f(y)-f(x)=(y-x)f'(c)$ for some $c\in(x,y)$,\\
so $\abs{f(y)-f(x)}\leq M(y-x)$. Symmetry handles $x>y$, and $x=y$ is trivial.\\
\hlt{Structural Commentary}: A derivative bound controls every secant slope.\\
\hlt{Key Insight}: Local slope bounds become global Lipschitz bounds through the Mean Value Theorem.
\end{proof}
\end{corollary}

\begin{theorem} \label{thm:derivative-sign-monotonicity} \hlt{Derivative Sign and Monotonicity}\\
Suppose $f$ is differentiable in $(a,b)$.
\begin{enumerate}[label=(\roman*)]
\setlength{\itemsep}{0pt}
\item If $f'(x)\geq0$ for all $x\in(a,b)$, then $f$ is monotonically increasing.
\item If $f'(x)=0$ for all $x\in(a,b)$, then $f$ is constant.
\item If $f'(x)\leq0$ for all $x\in(a,b)$, then $f$ is monotonically decreasing.
\end{enumerate}
\begin{proof} For $x_1<x_2$ in $(a,b)$, Corollary \ref{cor:mean-value-theorem} on $[x_1,x_2]$ gives $f(x_2)-f(x_1)=(x_2-x_1)f'(c)$ for some $c\in(x_1,x_2)$.
\begin{enumerate}[label=(\roman*)]
\setlength{\itemsep}{0pt}
\item If $f'\geq0$, then $f(x_2)-f(x_1)\geq0$.
\item If $f'=0$, then $f(x_2)-f(x_1)=0$.
\item If $f'\leq0$, then $f(x_2)-f(x_1)\leq0$.
\end{enumerate}
\hlt{Structural Commentary}: Monotonicity is a global reading of the mean-value identity.\\
\hlt{Key Insight}: The sign of every tangent slope controls every secant slope.
\end{proof}
\end{theorem}

\begin{theorem} \label{thm:derivatives-have-intermediate-value-property} \hlt{Derivatives Have the Intermediate Value Property}\\
If $f$ is differentiable on $[a,b]$ in the one-sided endpoint sense and $f'(a)<\lambda<f'(b)$,\\
Then $f'(x)=\lambda$ for some $x\in(a,b)$. The same conclusion holds when $f'(b)<\lambda<f'(a)$.
\begin{proof}
Set $g(t)=f(t)-\lambda t$. Since $g'(a)<0<g'(b)$, there are $t_1,t_2\in(a,b)$ with $g(t_1)<g(a)$ and $g(t_2)<g(b)$. Hence the minimum of $g$ on $[a,b]$ occurs at some $x\in(a,b)$. By Theorem \ref{thm:interior-extremum-derivative-zero}, $0=g'(x)=f'(x)-\lambda$.\\
Apply the same argument to $-f$ for the reversed case.\\
\hlt{Structural Commentary}: A derivative need not be continuous, but it cannot jump over values.\\
\hlt{Key Insight}: Tilt the graph by subtracting $\lambda t$ and force an interior extremum.
\end{proof}
\end{theorem}

\begin{figure}[H]
  \centering
  \begin{tikzpicture}[
  x=0.72cm,
  y=0.67cm,
  >=Stealth,
  every node/.style={font=\scriptsize},
  axis/.style={->, black!55, line width=0.45pt},
  curve/.style={black!85, line width=0.9pt},
  secant/.style={black!65, line width=0.8pt},
  guide/.style={black!22, dashed, line width=0.38pt},
  tick/.style={black!55, line width=0.45pt},
  point/.style={circle, fill=black, inner sep=1.15pt},
  hbrace/.style={<->, black!55, line width=0.45pt}
]
  \draw[axis] (-0.35,0) -- (8.45,0) node[right] {$t$};
  \draw[axis] (0,-0.25) -- (0,4.15);

  \coordinate (A) at (0.75,3.40);
  \coordinate (B) at (1.65,2.00);
  \coordinate (C) at (3.75,1.42);
  \coordinate (D) at (4.65,1.56);
  \coordinate (E) at (6.05,2.65);
  \coordinate (F) at (6.95,3.58);

  \draw[curve]
    (0.30,3.90) .. controls (0.55,3.65) and (0.62,3.52) .. (A)
    .. controls (1.00,2.72) and (1.28,2.18) .. (B)
    .. controls (2.25,1.08) and (3.12,1.05) .. (C)
    .. controls (4.05,1.35) and (4.32,1.48) .. (D)
    .. controls (5.15,1.72) and (5.55,2.18) .. (E)
    .. controls (6.35,3.00) and (6.55,3.28) .. (F)
    .. controls (7.35,3.82) and (7.70,3.72) .. (8.05,3.52);

  \foreach \p/\q in {A/B, C/D, E/F} {
    \draw[secant] (\p) -- (\q);
    \node[point] at (\p) {};
    \node[point] at (\q) {};
  }

  \foreach \x in {0.75,1.65,3.75,4.65,6.05,6.95} {
    \draw[guide] (\x,0) -- (\x,3.75);
    \draw[tick] (\x,0.09) -- (\x,-0.09);
  }

  \foreach \x/\lab in {0.75/{x_1},3.75/{x}} {
    \node[below] at (\x,-0.08) {$\lab$};
  }
  \node[below] at (6.95,-0.08) {$x_2$};

  \foreach \xl/\xr in {0.75/1.65,3.75/4.65,6.05/6.95} {
    \draw[hbrace] (\xl,-0.58) -- node[below] {$h$} (\xr,-0.58);
  }

  \draw[->, black!45] (2.15,0.52) -- (3.18,0.52);
  \draw[->, black!45] (5.05,0.52) -- (5.62,0.52);
\end{tikzpicture}

  \caption{Sliding fixed-width secants in Darboux's theorem.}
  \label{fig:sliding-secant-darboux-derivative}
\end{figure}

\begin{perspective} \label{per:sliding-secant-proof-of-darboux} \hlt{Sliding Secants}\\
For fixed small $h>0$, the secant-slope function $\sigma_h(x)\defeq \frac{f(x+h)-f(x)}{h}$ is continuous. Thus Theorem \ref{thm:intermediate-value-theorem} finds intermediate secant slopes, and Corollary \ref{cor:mean-value-theorem} converts them to derivative values.\\
\hlt{Key Insight}: Slide secants, then apply the Mean Value Theorem.
\end{perspective}

\begin{corollary} \label{cor:derivatives-no-simple-discontinuities} \hlt{Derivatives Have No Simple Discontinuities}\\
If $f$ is differentiable on $[a,b]$, then $f'$ cannot have a simple discontinuity on $[a,b]$.
\begin{proof}
Suppose a simple discontinuity occurs at an interior point $c$; endpoint cases are one-sided.\\
If $f'(c-)\neq f'(c+)$, choose $\lambda$ strictly between these limits, $\lambda\neq f'(c)$. For $x<c<y$ close to $c$, Theorem \ref{thm:derivatives-have-intermediate-value-property} gives $z\in(x,y)$ with $f'(z)=\lambda$, contradicting the one-sided limits. If $f'(c-)=f'(c+)=L\neq f'(c)$, choose $\lambda$ between $L$ and $f'(c)$ and apply the same theorem on a sufficiently small one-sided interval with endpoint $c$.\\
\hlt{Structural Commentary}: Darboux behavior rules out jump discontinuities for derivatives.\\
\hlt{Key Insight}: Derivatives may oscillate, but they do not jump.
\end{proof}
\end{corollary}

\begin{theorem} \label{thm:inverse-function-derivative-one-variable} \hlt{Derivative of an Inverse Function}\\
Let $f:I\to\R$ be one-to-one on an open interval $I$, let $x_0\in I$, and put $y_0=f(x_0)$. Suppose $f$ is differentiable at $x_0$, $f'(x_0)\neq0$, and $g=f^{-1}:f(I)\to I$ is continuous at $y_0$. Then $g$ is differentiable at $y_0$ and $g'(y_0)=\frac{1}{f'(x_0)}$.
\begin{proof}
For $y\in f(I)$ with $y\neq y_0$, write $x=g(y)$. Then $x\neq x_0$ and $\frac{g(y)-g(y_0)}{y-y_0} =\frac{x-x_0}{f(x)-f(x_0)} =\left(\frac{f(x)-f(x_0)}{x-x_0}\right)^{-1}$.\\
As $y\to y_0$, the assumed continuity of $g$ gives $x=g(y)\to x_0$.\\
The difference quotient for $f$ tends to $f'(x_0)\neq0$, so the reciprocal tends to $1/f'(x_0)$ by Theorem \ref{thm:function-limit-algebra}.\\
\hlt{Structural Commentary}: Matched inverse quotients are reciprocal.\\
\hlt{Key Insight}: Continuity of the inverse converts the limit $y\to y_0$ back into $x\to x_0$.
\end{proof}
\end{theorem}

\begin{corollary} \label{cor:one-dimensional-inverse-function-theorem} \hlt{One-Dimensional Inverse Function Theorem}\\
If $f:(a,b)\to(c,d)$ is a differentiable bijection and $f'\neq0$,\\
then $f$ is a homeomorphism and $f^{-1}$ is differentiable, with $(f^{-1})'(y)=\frac{1}{f'(f^{-1}(y))}$.
\begin{proof}
By Theorem \ref{thm:derivatives-have-intermediate-value-property}, $f'$ has constant sign. Corollary \ref{cor:mean-value-theorem} makes $f$ strictly monotone, hence its continuous interval inverse is continuous. Apply Theorem \ref{thm:inverse-function-derivative-one-variable} pointwise.\\
\hlt{Structural Commentary}: Darboux gives constant sign; the inverse derivative formula is local.\\
\hlt{Key Insight}: Nonvanishing derivative is the one-dimensional obstruction to differentiating the inverse.
\end{proof}
\end{corollary}

\begin{theorem} \label{thm:lhospital-one-sided} \hlt{L'Hospital's Rule}\\
Let $f,g$ be differentiable on $(a,b)$, $g'\neq0$, and $-\infty\leq a<b\leq+\infty$.\\
If $\frac{f'}{g'}\to A\in\overline{\R}$ as $x\to a$, and either $f,g\to0$ or $g\to+\infty$, then $\frac{f}{g}\to A$ as $x\to a$.\\
The same holds as $x\to b$, and for $g\to-\infty$.
\begin{proof}
For $a<x<y<b$, Theorem \ref{thm:cauchy-mean-value-theorem} gives $t\in(x,y)$ with $\frac{f(x)-f(y)}{g(x)-g(y)}=\frac{f'(t)}{g'(t)}$.\\
If $f,g\to0$, let $x\to a$ with $y$ fixed. If $g\to+\infty$, use $\frac{f(x)}{g(x)} =\frac{f(x)-f(y)}{g(x)-g(y)}\,\frac{g(x)-g(y)}{g(x)}+\frac{f(y)}{g(x)}$ and let $x\to a$.\\
In both cases, bounds on $f'/g'$ near $a$ transfer to $f/g$, forcing the limit.\\
\hlt{Structural Commentary}: CMVT transfers derivative bounds to quotient bounds.\\
\hlt{Key Insight}: Derivative quotients govern the original indeterminate quotient.
\end{proof}
\end{theorem}

\begin{figure}[H]
  \centering
  \begin{tikzpicture}[
  x=0.97cm,
  y=0.72cm,
  >=Stealth,
  every node/.style={font=\scriptsize},
  axis/.style={->, gray!70, line width=0.45pt},
  curve/.style={blue!70!black, line width=0.9pt},
  guide/.style={gray!50, dashed, line width=0.4pt},
  labelbox/.style={fill=white, fill opacity=0.9, text opacity=1, inner sep=1.2pt}
]
  \draw[axis] (-2.65,0) -- (2.75,0) node[right] {$x$};
  \draw[axis] (0,-2.55) -- (0,2.75) node[above] {$y$};

  \begin{scope}
    \clip (-2.55,-2.45) rectangle (2.55,2.45);
    \draw[curve, domain=-2.25:-0.36, samples=150]
      plot (\x,{0.45/(\x*\x*\x)});
    \draw[curve, domain=0.36:2.25, samples=150]
      plot (\x,{0.45/(\x*\x*\x)});
  \end{scope}

  \draw[curve, ->, domain=-0.78:-0.61, samples=30]
    plot (\x,{0.45/(\x*\x*\x)});
  \draw[curve, ->, domain=0.78:0.61, samples=30]
    plot (\x,{0.45/(\x*\x*\x)});
  \node[labelbox, align=center] at (-1.55,-1.55) {$x\to0^-$\\[-1pt]$f(x)\to-\infty$};
  \node[labelbox, align=center] at (1.55,1.55) {$x\to0^+$\\[-1pt]$f(x)\to+\infty$};
  \node[below left] at (0,0) {$0$};
\end{tikzpicture}

  \caption{For $f(x)=x^{-3}$, the one-sided limits at $0$ are opposite infinities.}
  \label{fig:reciprocal-cubic-one-sided-limits}
\end{figure}

\begin{definition} \label{def:higher-derivatives} \hlt{Higher Derivatives}\\
If $f$ has derivative $f'$ on an interval and $f'$ is differentiable, then $(f')'$ is denoted $f''$ and called the \hlt{second derivative}. Continuing recursively gives $f^{(n)}$, the derivative of order $n$.
\end{definition}

\begin{definition} \label{def:smoothness-classes-one-variable} \hlt{Smoothness Classes}\\
On an interval, $C^0$ means continuous, $C^r$ means $f^{(r)}$ exists and continuous, $C^\infty$ means $C^r$ for every finite $r$.\\
Analytic functions (Definition \ref{def:real-analytic-at-point}) form the class $C^\omega$.
\end{definition}

\begin{perspective} \label{per:regularity-hierarchy-one-variable} \hlt{Regularity Hierarchy}\\
The inclusions are $C^\omega\subseteq C^\infty\subseteq \cdots \subseteq C^2\subseteq C^1\subseteq C^0$.\\
The inclusions follow from Theorem \ref{thm:power-series-uniform-differentiation} and Definition \ref{def:smoothness-classes-one-variable}; they are generally proper.\\
\hlt{Key Insight}: Higher regularity is extra structure, not just a different notation for continuity.
\end{perspective}

\begin{theorem} \label{thm:higher-order-leibniz-product-rule} \hlt{Higher-Order Leibniz Product Rule}\\
If $f,g:(a,b)\to\R$ are $r$ times differentiable, $r\geq1$, then $(fg)^{(r)}(x)=\sum\limits_{k=0}^{r}\binom{r}{k}f^{(k)}(x)g^{(r-k)}(x)$.
\begin{proof}
Induct on $r$. The case $r=1$ is Theorem \ref{thm:derivative-algebra-rules}.\\
Differentiating the order-$r$ formula gives $(fg)^{(r+1)} =\sum\limits_{k=0}^{r}\binom{r}{k}f^{(k+1)}g^{(r-k)} + \sum\limits_{k=0}^{r}\binom{r}{k}f^{(k)}g^{(r-k+1)}$.\\
Reindex and use Pascal's identity.\\
\hlt{Structural Commentary}: Each derivative chooses whether it lands on $f$ or on $g$.\\
\hlt{Key Insight}: The binomial coefficients count those choices.
\end{proof}
\end{theorem}

\begin{theorem} \label{thm:taylor-theorem-lagrange-remainder} \hlt{Taylor's Theorem with Lagrange Remainder}\\
Let $f:[a,b]\to\R$, let $n\geq1$, suppose $f^{(n-1)}$ is continuous on $[a,b]$, and suppose $f^{(n)}$ exists on $(a,b)$.\\
For distinct $\alpha,\beta\in[a,b]$, set $P(t)\defeq\sum\limits_{k=0}^{n-1}\frac{f^{(k)}(\alpha)}{k!}(t-\alpha)^k$.\\
Then for some $x$ between $\alpha$ and $\beta$, $f(\beta)=P(\beta)+\frac{f^{(n)}(x)}{n!}(\beta-\alpha)^n$.
\begin{proof}
Choose $M$ so that $f(\beta)=P(\beta)+M(\beta-\alpha)^n$, and set $G(t)\defeq f(t)-P(t)-M(t-\alpha)^n$.\\
Then $G(\alpha)=G'(\alpha)=\cdots=G^{(n-1)}(\alpha)=G(\beta)=0$. Repeated applications of Corollary \ref{cor:mean-value-theorem} give $G^{(n)}(x)=0$ for some $x$ between $\alpha$ and $\beta$. Since $G^{(n)}(x)=f^{(n)}(x)-n!M$, the formula follows.\\
\hlt{Structural Commentary}: The auxiliary error vanishes at $\beta$ and to order $n-1$ at $\alpha$.\\
\hlt{Key Insight}: Repeated Rolle-style arguments turn endpoint cancellation into a precise remainder.
\end{proof}
\end{theorem}

\begin{example} \label{ex:smooth-flat-nonanalytic-function} \hlt{Smooth but Not Analytic}\\
Define $e(x)=
\begin{cases}
e^{-1/x},&x>0,\\
0,&x\leq0.
\end{cases}$.
By repeated L'Hospital, Theorem \ref{thm:lhospital-one-sided}, $e\in C^\infty$ and $e^{(r)}(0)=0$ for all $r$.\\
Its Taylor series at $0$ is zero, but $e(x)>0$ for $x>0$, so $e$ is not analytic at $0$ (Definition \ref{def:real-analytic-at-point}).\\
\hlt{Key Insight}: A function can be infinitely tangent to the $x$-axis without being locally equal to its Taylor series.
\end{example}

\begin{method} \label{meth:newton-and-secant-iterations} \hlt{Newton and Secant Iterations}\\
To solve $f(x)=0$:
\begin{enumerate}[label=(\roman*)]
\setlength{\itemsep}{0pt}
\item
Newton uses tangent-line zero: $x_n=x_{n-1}-\frac{f(x_{n-1})}{f'(x_{n-1})}$, $n\geq1$ when $f'(x_{n-1})\neq0$.
\item
Secant uses two-point slope: $x_n=x_{n-1}-\frac{f(x_{n-1})(x_{n-2}-x_{n-1})}{f(x_{n-2})-f(x_{n-1})}$, $n\geq2$ when the denominator is nonzero.
\end{enumerate}
These define iterates only; local convergence is proved with Taylor's theorem (Theorem \ref{thm:taylor-theorem-lagrange-remainder}).
\end{method}

\begin{figure}[H]
  \centering
  \begin{tikzpicture}[
  x=0.72cm,
  y=0.72cm,
  >=Stealth,
  every node/.style={font=\scriptsize},
  axis/.style={->, black!55, line width=0.45pt},
  curve/.style={blue!70!black, line width=0.9pt},
  tangent/.style={green!45!black, line width=0.75pt},
  secant/.style={orange!85!black, line width=0.75pt},
  guide/.style={black!28, dashed, line width=0.38pt},
  tick/.style={black!55, line width=0.45pt},
  point/.style={circle, fill=black!80, inner sep=1.35pt},
  labelbox/.style={fill=white, fill opacity=0.92, text opacity=1, inner sep=1.1pt}
]
  \begin{scope}
    \draw[axis] (-0.25,0) -- (6.85,0) node[right] {$x$};
    \draw[axis] (0,-1.05) -- (0,1.55);

    \draw[curve, domain=0.35:6.45, samples=110]
      plot (\x,{0.15*(\x-3.65)*(\x-3.65)-0.64});

    \coordinate (P) at (6.20,0.335);
    \coordinate (R) at (5.762,0);

    \draw[tangent] (5.00,-0.583) -- (6.55,0.603);
    \draw[guide] (P) -- (6.20,0);
    \draw[guide] (R) -- (5.762,-0.50);

    \node[point] at (P) {};
    \fill[green!45!black] (R) circle (1.35pt);

    \draw[tick] (5.762,0.08) -- (5.762,-0.08);

    \node[labelbox] at ($(R)+(0,-0.55)$)
      {\textcolor{green!45!black}{$x_{n+1}$}};
    \node[labelbox, anchor=south west] at ($(P)+(0.05,0.04)$) {$x_n$};
    \node[anchor=west] at (0.18,1.35) {\textcolor{green!45!black}{Newton}};
  \end{scope}

  \begin{scope}[xshift=7.15cm]
    \draw[axis] (-0.25,0) -- (6.85,0) node[right] {$x$};
    \draw[axis] (0,-1.05) -- (0,1.55);

    \draw[curve, domain=0.35:6.45, samples=110]
      plot (\x,{0.15*(\x-3.65)*(\x-3.65)-0.64});

    \coordinate (A) at (4.25,-0.586);
    \coordinate (B) at (6.20,0.335);
    \coordinate (R) at (5.490,0);

    \draw[secant] (3.70,-0.846) -- (6.55,0.500);
    \draw[guide] (A) -- (4.25,0);
    \draw[guide] (B) -- (6.20,0);
    \draw[guide] (R) -- (5.490,-0.50);

    \node[point] at (A) {};
    \node[point] at (B) {};
    \fill[orange!85!black] (R) circle (1.35pt);

    \draw[tick] (5.490,0.08) -- (5.490,-0.08);

    \node[labelbox] at ($(R)+(0,-0.55)$)
      {\textcolor{orange!85!black}{$x_{n+1}$}};
    \node[labelbox, anchor=north east] at ($(A)+(-0.08,-0.13)$) {$x_{n-1}$};
    \node[labelbox, anchor=south west] at ($(B)+(0.05,0.04)$) {$x_n$};
    \node[anchor=west] at (0.18,1.35) {\textcolor{orange!85!black}{Secant}};
  \end{scope}
\end{tikzpicture}

  \caption{\textcolor{green!45!black}{Newton}: tangent-line zero. \textcolor{orange!85!black}{Secant}: two-point line zero.}
  \label{fig:newton-secant-iterations}
\end{figure}

\begin{definition} \label{def:vector-valued-derivative-one-variable} \hlt{Vector-Valued Derivative}\\
For $\mathbf{f}:[a,b]\to\R^k$, set $\mathbf{f}'(x)=\mathbf{L}$ when $\lim\limits_{t\to x}\norm{\frac{\mathbf{f}(t)-\mathbf{f}(x)}{t-x}-\mathbf{L}}_2=0$.\\
For $\mathbf{f}=(f_1,\ldots,f_k)$, this is equivalent to componentwise differentiability, with $\mathbf{f}'(x)=(f_1'(x),\ldots,f_k'(x))$.
\end{definition}

\begin{theorem} \label{thm:vector-valued-mean-value-inequality} \hlt{Vector-Valued Mean Value Inequality}\\
If $\mathbf{f}:[a,b]\to\R^k$ is continuous on $[a,b]$ and differentiable on $(a,b)$.\\
Then for some $x\in(a,b)$, $\norm{\mathbf{f}(b)-\mathbf{f}(a)}_2\leq (b-a)\norm{\mathbf{f}'(x)}_2$.
\begin{proof}
Set $\mathbf{z}\defeq\mathbf{f}(b)-\mathbf{f}(a)$ and $\varphi(t)\defeq\mathbf{z}\cdot\mathbf{f}(t)$.\\
Corollary \ref{cor:mean-value-theorem} gives $\varphi(b)-\varphi(a)=(b-a)\varphi'(x)=(b-a)\mathbf{z}\cdot\mathbf{f}'(x)$ for some $x\in(a,b)$. Since $\varphi(b)-\varphi(a)=\norm{\mathbf{z}}_2^2$, Theorem \ref{thm:finite-schwarz-inequality-complex-scalars} gives $\norm{\mathbf{z}}_2^2\leq(b-a)\norm{\mathbf{z}}_2\norm{\mathbf{f}'(x)}_2$, and the result follows by cancellation, with $\mathbf{z}=0$ trivial.\\
\hlt{Structural Commentary}: A scalar projection yields the vector norm bound.\\
\hlt{Key Insight}: Project onto the displacement vector, then use Schwarz to return to norms.
\end{proof}
\end{theorem}

\begin{method} \label{meth:one-variable-differentiation-no-overlap-source-todo}\label{meth:one-variable-differentiation-proof-patterns} \hlt{Differentiation Proof Patterns}
\begin{enumerate}[label=(\roman*)]
\setlength{\itemsep}{0pt}
\item \hlt{Domain}: require $x_0$ to be a limit point of $X$; isolated points have no derivative by Perspective \ref{per:one-variable-domain-of-differentiability}.
\item \hlt{Definition}: factor $f(t)-f(x)=(t-x)Q_x(t)$ and take $Q_x(t)\to f'(x)$.
\item \hlt{Approximation}: use Proposition \ref{prop:one-variable-first-order-approximation} to replace quotient limits by first-order error bounds.
\item \hlt{Algebra}: prove quotient identities first; use continuity to pass to limits.
\item \hlt{Mean values}: choose $h$ with equal endpoint values and apply Theorem \ref{thm:interior-extremum-derivative-zero}.
\item \hlt{Taylor}: set $G=f-P-M(t-\alpha)^n$ and iterate the mean value theorem.
\end{enumerate}
\hlt{Key Insight}: Difference quotients and auxiliary functions drive one-variable differentiation.
\end{method}
